\begin{table}[htbp]
\centering
\caption{P1-series results and interpretation boundaries. The five seeds are algorithmic repetitions; the number of physical trajectories remains three.}
\label{tab:p1-series-summary}
\small
\begin{tabularx}{\textwidth}{@{}l>{\raggedright\arraybackslash}X>{\raggedright\arraybackslash}X@{}}
\toprule
Experiment & Numerical result & Domain of validity \\
\midrule
P1 &
Standard-CHNN joint-field, rollout, and linearization errors are \(0.1417\%\)--\(0.2444\%\), \(0.2131\%\)--\(1.5237\%\), and \(0.3125\%\)--\(0.3286\%\). &
Restricted to a three-trajectory, derivative-only small-sample comparison. \\
P1.5 &
Residual-over-physics-only mean reductions are \(90.37\%\)--\(97.29\%\) for the field and \(53.40\%\)--\(95.14\%\) for rollout; the CHNN-over-gray-Real mean field reduction is \(-2.195\%\). &
Conditional and more expensive; no stable one-step-field advantage. \\
P1.6 &
All \(30/30\) constrained models pass the equilibrium gates exactly, while the mean Real/CHNN joint-field errors worsen to \(2.872\times10^{-2}\)/\(2.623\times10^{-2}\). &
Constraint satisfaction does not imply improved accuracy. \\
\bottomrule
\end{tabularx}
\end{table}
